\section{Two raw-kernel realizations of the general theorem}\label{sec:realizations}
The first example has classical correlated reports and a polynomial test closure. The second has noncommuting quantum effects and action-dependent projective supports. Both verify \eqref{eq:envelopes} directly and hence instantiate the same Theorem~\ref{thm:central}; neither is used to prove it.

\begin{theorem}[Controlled classical observation]\label{thm:classical}
Let the physical state space be $[-1/2,1/2]^d$, and let each action have a centered compact probability $\nu_a$ on that box. Prepare $x_0=0$. For unknown fixed $\theta=(\theta_a)_a\in[-\eta,\eta]^{\calA}$, $0\le\eta<1$, define
\begin{equation}\label{eq:classical-kernel}
 K_a^\theta(dx',dy\mid x)
 =\left(1+\frac{4\theta_a}{d}\ip{x}{y}\right)\delta_y(dx')\nu_a(dy).
\end{equation}
After a stop, audit the coordinates independently with values in $\{-1,0,1\}$ and probabilities
\[
 \Pp(Z_i=\pm1\mid x)=\frac{x_i^2+1/2\pm x_i}{2},\qquad
 \Pp(Z_i=0\mid x)=\frac12-x_i^2.
\]
Then the task mean is $X_t=Y_t$, its conditional variance is $B=d/2$, and Theorem~\ref{thm:central}\textup{(iii)} holds uniformly in the unknown parameter with $\underline b_a=1-\eta$, $\overline b_a=1+\eta$. If all $\theta_a=0$, its exact parts apply with the different laws $\nu_a$. Every regular, singular or atomic choice of these acquired laws satisfies the corresponding conclusions whenever their separately stated geometric hypotheses hold.
\end{theorem}
\begin{proof}
The likelihood is between $1-\eta$ and $1+\eta$, since $|\ip{x}{y}|\le d/4$. Centering of $\nu_a$ gives normalization. The actual first-call law from $x_0=0$ is $\nu_a$, for every parameter. The report is the new state, and neither a rejected report nor the old state is reset to the preparation. The displayed row inequalities hold at every hidden state, hence at every conditional mixture induced by a history and any common policy. The audit probabilities are nonnegative on the box, sum to one, and give $\E Z_i=x_i$, $\E Z_i^2=x_i^2+1/2$. This proves the mean and variance assertions.

These formulas arise from a finite affine executable-test representation rather than an assumed Fisher geometry. Let $\phi(x)$ be the vector of all monomials $x^\alpha$ with $\alpha_i\in\{0,1,2\}$, including one, and let its positive cone be generated by $\phi(x)$ on the physical box. The raw report instrument acts by
\[
 L_a^\theta(dy)z
 =\left(z_0+\frac{4\theta_a}{d}\sum_i y_i z_{e_i}\right)\phi(y)\nu_a(dy).
\]
It is positive on that cone and normalized on its base. The full protected audit probabilities are linear combinations of these monomials. Report updates, future actions and audit tests therefore lie in a genuine raw positive realization. No multiplication closure beyond this explicitly verified family is presumed.

For two incoming physical states the report laws have total variation at most $\eta\norm{x-x'}/\sqrt d$ by the displayed affine density. The same-report current-task update is $y$ for both states, so it has zero discrepancy. The global covers and small-ball constants, when required, are exactly those of the actually acquired $\nu_a$. The program in the central theorem depends on these known measures, not on $\theta$.
\end{proof}

Centering and a fixed affine rescaling embed the atomic--uniform law in Theorem~\ref{thm:singular} into this box; squared distortions acquire the square of the scale factor. One may take centered Cantor measures, smooth factors, or mixed dimensions for different actions. Their preparation and repeated acquisition costs remain those of \eqref{eq:classical-kernel}.

\begin{theorem}[Noncommuting action-dependent quantum supports]\label{thm:quantum}
Let $q=2^k$, $k\ge1$, and for each action choose $r_a=2^{j_a}$ with $1\le j_a\le k$, an isometry $V_a:\mathbb C^{r_a}\to\mathbb C^q$, and a trace-preserving positive channel $C_a$ with
\[
 C_a(I_q/q)=I_{r_a}/r_a.
\]
For example $C_a(\rho)$ may be a partial trace of a unitary conjugation on a factorization $\mathbb C^q=\mathbb C^{r_a}\otimes\mathbb C^{q/r_a}$. Set
\[
 T_a(\rho)=(1-\eta_a)C_a(\rho)+\eta_a\tr(\rho)I_{r_a}/r_a,
 \qquad \eta_a\in[\eta_0,1],\quad\eta_0>0,
\]
and define the raw instrument, relative to Haar probability $\sigma_{r_a}$ on projective pure states, by
\begin{equation}\label{eq:quantum-kernel}
 L_a(d\psi)\rho
 =r_a\tr(P_\psi T_a(\rho))\,V_aP_\psi V_a^*\,\sigma_{r_a}(d\psi).
\end{equation}
Prepare $I_q/q$. The channels and $\eta_a$ may range over a known family with a fixed unknown member; the isometries and the current-state readout are known. Use the protected covariant audit on $\mathbb C^q$ with density $q\tr(\rho P_\zeta)$ and readout $Z=(q+1)P_\zeta-I_q$.

With Hilbert--Schmidt squared loss, the task mean is $X=V_aP_\psi V_a^*$, $B=q^2+q-2$, and the central envelope theorem applies with $\underline b_a=\eta_0$, $\overline b_a=r_a$. Its actual reference support has dimension $s_a=2(r_a-1)$ and satisfies
\begin{equation}\label{eq:projective-mass}
 \nu_a\{P:\norm{P-P_0}_\HS\le u\}
 =(u^2/2)^{r_a-1}\qquad(0\le u\le\sqrt2).
\end{equation}
The channels need not commute, and no common eigenbasis is used. When $\eta_a=1$ for all actions the task is regenerative and the exact action profile applies.
\end{theorem}
\begin{proof}
The effect density in \eqref{eq:quantum-kernel} is between $\eta_0$ and $r_a$ on normalized positive inputs. Its integral is one because $\int P_\psi\,d\sigma_{r_a}=I_{r_a}/r_a$. The output normalized state is precisely the reported $V_aP_\psi V_a^*$, and the report law from the specified preparation is Haar for each action. Mixtures of unknown past states preserve the density bounds. These prove all conditional-envelope quantifiers, including histories after rejections.

For the audit, unitary invariance gives
\[
 \int P_\zeta\otimes P_\zeta\,d\sigma_q
       =\frac{I+\mathrm{Swap}}{q(q+1)}.
\]
Indeed the integrand is supported on the symmetric subspace, its integral is scalar there by invariance, and its trace is one; the symmetric dimension is $q(q+1)/2$. Partial tracing the identity proves $\E[Z\mid\rho]=\rho$. Also $\norm Z_\HS^2=q^2+q-1$ for every audit outcome. Since every stopped state from \eqref{eq:quantum-kernel} is pure, its conditional variance is $q^2+q-2$, independent of the action and parameter.

This is a raw finite-dimensional positive instrument on Hermitian matrices, with normalization $\tr$. Future instruments and the audit are positive linear operations before normalization. No quantum system is available as free observer memory: only the classical labels in Section~\ref{sec:experiment} are retained by the observer. For normalized inputs $\rho,\rho'$, report total variation is at most $\frac12\norm{\rho-\rho'}_1$, by positivity and trace preservation of the measurement channel. The current-task outputs agree for a common action and report. Covers and small-ball masses are obtained directly from \eqref{eq:projective-mass}. To verify that formula, choose $P_0=e_1e_1^*$; the random variable $|\psi_1|^2$ has density $(r_a-1)(1-t)^{r_a-2}$, as follows by normalizing independent complex Gaussian coordinates. Since $\norm{P_\psi-P_0}_\HS^2=2(1-|\psi_1|^2)$, integration yields the displayed mass. These steps prove the same general theorem independently of the classical realization.
\end{proof}

A marked convex mixture of \eqref{eq:quantum-kernel} with the positive operation $\rho\mapsto\tr(\rho)P_0$ produces an actual atom plus a continuous projective component. Choose $P_0$ outside the selected proper projective subspace when separation is desired. Both outputs remain pure, so the same audit variance applies. For $\eta_a=1$ the mixed law is regenerative, and Theorem~\ref{thm:singular} gives its mass-dependent frontier. For $\eta_a<1$ the conditional-envelope conclusions apply. This verifies a singular stratum in a noncommuting raw experiment; an arbitrary rank-changing filter is not being identified with it.
